\section{Risks and critical reflection, in full}\label{sec:I}

\S\ref{sec:13} states each finding, its verdict and its bound. This appendix carries the mechanism behind
each one. The subsection numbering follows \S\ref{sec:13} so a reader can move between them.

\subsection{Four functional holes, two of which are now closed (\texorpdfstring{\S\ref{sec:13.1}}{13.1} in full)}\label{sec:I.1}

\textbf{A damper stuck \emph{wide open} is invisible at the deployed setting.} The statistic is
two-sided: the lower arm is the exact mirror of the upper one, $S_\text{low} = \max(0, S_\text{low}
- r - k)$, sharing $k$ and $h$ so that symmetry cost no new tuning parameter. The deployed default
is the \emph{guarded} lower arm, driven by the measurement against the trajectory of an empty room
clamped at the 420 ppm outdoor concentration, and it is what Section \ref{sec:11.2}'s headline figures are
measured on. Two experiments deliberately run the arm \emph{off} beside it, as a control rather
than as a configuration: 180 of T-34b's 540 observed rooms and one of T-26's four regimes. That guard is
what catches the low-side fault with a safety consequence: a sensor drifting downward makes the
inner loop believe the room is cleaner than it is and \emph{under-ventilate an occupied room}, and
it betrays itself by reading below the outdoor concentration, which no amount of ventilation can
produce. What the guard cannot catch is the other low-side fault. ``Everybody left the room''
explains a room that stopped getting stuffier at least as well as ``the damper ignored the close
command'' does, and nothing inside a ten-minute judging window separates them, so a damper stuck
wide open stays \texttt{HEALTHY}. The unguarded setting sees both and costs more false confirmations
than R-3 allows, which the monitor's own fault-free-week test measures. That is the honest form of
this hole: an observer with no occupancy signal \emph{cannot} have both, and the choice was made in
R-3's favour.

\textbf{\texttt{serviced()} had no transport, so a latched false positive could not be cleared in
deployment. And now it has one.} \texttt{FaultCase.serviced()} is the only exit from a latched
case, and for most of this project's life nothing reached it: no MQTT topic, no REST endpoint, no
dashboard action, only the unit tests and the harness. A false latch was \textbf{permanent} and the
technician who repaired the damper had no way to tell the system so. That hole is closed.
\texttt{ServiceReported} travels on \texttt{service/\{zone\}}, \texttt{fault-monitor} subscribes to
it through an adapter that validates the report against the open case, and the dashboard exposes
\texttt{POST /api/zones/\allowbreak \{zone\}/cases/\allowbreak \{caseId\}/serviced}, answering 202 when it is accepted, 404
for a room or case that does not exist, 409 when that case is not latched, and 503 when the broker
is unreachable. The whole
\texttt{CONFIRMED}~→~\texttt{COMPENSATING}~→~\texttt{SERVICED}~→~\texttt{WARMUP}~→~\texttt{HEALTHY}
lifecycle is executable and tested end to end, including that a reported repair unwinds the
compensation and hands the five healthy rooms back to the inner loop. \textbf{This matters most to
the latching trade-off below}, whose whole argument is that a false positive costs only one
technician visit: an argument that was being made against a system in which it cost that visit
\emph{and} a permanently wrong work order. It is now true as stated. One thing is deliberately
absent: there is no button on the HTML status page, because that page is server-rendered with a
meta refresh and no JavaScript, so a form post would land the operator on a JSON body. The endpoint
is the operator path.

\textbf{Blocked filters are plant-wide by construction, so the classifier's hardest discrimination
is never tested.} The statement ``the three signatures may collapse into two'' is too soft. In
\emph{this} simulator they cannot collapse: a blocked filter is plant-wide by construction, so it
starves all six rooms and a stuck damper never does, and the common-mode count separates them
trivially. The ambiguity that would really arise in a building (a partial obstruction in one duct
branch, which reads as both a partly stuck damper and a partly blocked filter) \textbf{is not in
the fault model at all}. The enumeration cannot express it. So a clean confusion matrix from the
classifier would be weak evidence rather than strong, and reporting one without this paragraph
would be misleading. The risk register forbade inventing a fourth \emph{signal} to rescue a
three-way taxonomy; it never forbade a fourth \emph{fault}, and adding the ambiguous one is what
would have made the classifier's evaluation worth reading.

\textbf{And worse than untested: the \texttt{BLOCKED\_FILTER} branch of the classifier is
unreachable with this physics.} This is the finding I least expected and the one I would lead with
in the oral. A unit test pins it deliberately: it injects a blocked filter and asserts the
classification is \texttt{UNCLASSIFIED}, with a message saying that if it ever starts classifying,
the test's finding is out of date. Two structural reasons, and neither is a bug to be patched:

\begin{itemize}
\tightlist
\item \textbf{The split only redistributes a fixed total.} $Q_i = Q_\text{fan}\,a_i/\sum a_j$ means
  no damper command can move all six residuals the same way; a blockage shrinks $Q_\text{fan}$, but
  the observer judges \emph{step responses to commands}, and no command it watches produces the
  plant-wide common mode the first branch of the tree requires.
\item \textbf{The generation estimate absorbs the blockage exactly.} From Section \ref{sec:7.2},
  $G_\text{est} = Q\,(C_\text{measured} - C_\text{out})$ is learned from quasi-steady state. With a
  blockage, the true flow is low by $(1-\beta)$ and the believed flow used in the prediction is high
  by the same factor, so $G_\text{est}$ comes out high by precisely $1/(1-\beta)$, and the two
  errors cancel in the predicted steady state. Only the \emph{time constant} differs, and the
  residual the CUSUM sees never develops the bias it needs. The observer's defence against not
  knowing occupancy is, in other words, also a defence against noticing a blocked filter.
\end{itemize}

What saves the outcome from being useless is that the classifier reports its evidence even when it
declines to name a cause, so the work order still states the measured plant power excess and that the
residual is not common mode: a technician is sent to the air handler with the right number, by a
system that admits it cannot prove which fault it is. That is the correct behaviour for an
unclassifiable case and it is not the same thing as a working three-way classifier. \textbf{The
honest headline is that of the mechanisms this project set out to name, exactly one is classified:
\texttt{STUCK\_DAMPER}.} \texttt{BLOCKED\_FILTER} can be pointed at but never confirmed, for the
structural reason above. And \texttt{SENSOR\_DRIFT}, which this paragraph used to call classified
``once a night of overnight-floor history exists'', is \textbf{never emitted at all}, in 0 of 76
runs from a drift ramp and 0 of 76 from a step offset (§\ref{sec:7.3}, §\ref{sec:11.5}). Its leaf wants an informative
overnight floor and a matched step response in the same instant, and the two injectable sensor
mechanisms supply one each and never both, because the classification is computed at the confirmation
while one of its inputs needs a day of the fault.

\subsection{The air handler's rating was corrected on measurement (\texorpdfstring{\S\ref{sec:13.3}}{13.3} in full)}\label{sec:I.2}

Run that plant at the normal 0.7 and it delivers 294 l/s against a 411.4 l/s design total, 71 \% of
design flow, and the consequence was not subtle: with \textbf{nothing broken} no room could be held
below 1000 ppm at design occupancy and the monitor confirmed a fault in all six, because a uniformly
starved building looks exactly like a faulty one. Sizing instead from $\Sigma$ design $\div\ 0.7$
separates the two jobs the rating was doing at once: the plant covers its design condition at the
speed it actually runs, and the margin to $s=1$ is reserved for compensation. The fix is not free,
and this is the honest cost: \textbf{margin masks faults}. At the corrected rating a filter blocking
35 \% no longer starves the rooms enough for the monitor to confirm it, where at the undersized
rating it did. Detection sensitivity was being bought with occupant comfort, which is the wrong
currency, so the trade is made deliberately in comfort's favour and its cost is stated here rather
than hidden. The correction is now measured rather than argued: at the corrected rating the
no-fault arm of T-26 holds every occupied minute at or below 1000 ppm at design occupancy, on all
five seeds, in every regime (§\ref{sec:11.2}). That is the figure the earlier sizing scored 0.000 on.


\subsection{The CUSUM margin, and why $h$ was never the knob (\texorpdfstring{\S\ref{sec:13.4}}{13.4} in full)}\label{sec:I.3}


The worst healthy seed peaked at 149.1 ppm against $h$ = 150 ppm at $\sigma$ = 30 ppm, over three
hours in one room: a \textbf{0.9 ppm margin}, with the next seed at 91.6 ppm, so it is a tail draw
and not a design property. R-3 asks for at most one false alarm per simulated week \emph{across six
rooms}: the same statistic over far more samples, six times over. The mitigation is persistence
rather than a higher limit: requiring the statistic to hold above $h$ for two consecutive cycles
removes a tail draw, which the next sample erases, while a stuck damper clears $h$ every cycle.
Raising $h$ instead trades R-3 against R-1, and the sweep in §\ref{sec:11.1} prices that trade: 8σ costs about
a minute of delay, and 4σ already gave one false alarm in twenty short runs. Siegmund's
ARL$_0 \approx 1.1\times10^{5}$ samples $\approx$ 79 days is analytic, on i.i.d. noise which the
production residual is not. And the measurement contradicts it. Over 20 fault-free simulated weeks
of all six rooms the monitor raises \textbf{0.55 confirmed cases per week against R-3's limit of
one}, with \textbf{16 of 20 weeks clean} and a worst week of 6, against 1.30 per week and 9 clean
weeks before ADR-0010, while on the same tuning R-1's stepped p95 is 6 min against a 15 min budget
(§\ref{sec:11.2}). So the trade this section prices came out on R-3's side of the line, but it was not the
binding trade, and that is the correction this section needs. \textbf{$h$ was never the knob.}
ADR-0010 attributed the failure, and the defect was in the \emph{reference} the limit is applied to
rather than in the limit: both plausibility guards were being corrected 30 \% of the way onto every
reading, so each CUSUM arm was accumulating a filter error instead of a residual, which is what put
the statistic into the very i.i.d. regime the analytic run length assumes and is why the analytic
figure was contradicted so badly. Correcting the lower arm's half of it halved the rate with
$k$ = 30 ppm and $h$ = 150 ppm \emph{untouched}. The response taken is therefore unchanged and
better supported than when it was taken: $k$ and $h$ stay where the spike put them rather than
being retuned until R-3 passes, because retuning a limit on the experiment that failed it is the
same methodological error §\ref{sec:13.10} already names, and the measurement now says the limit was not the
defect. What R-3 still fails on is the worst week, which is seed 2 at six confirmations, all on the
upper arm; the upper arm's half of the same defect is measured, meets R-3 outright, and was
declined because it costs three other requirements (§\ref{sec:11.2}, ADR-0010).

\subsection{R-7 at design occupancy, and the two bounds on that finding (\texorpdfstring{\S\ref{sec:13.5}}{13.5} in full)}\label{sec:I.4}

Until this revision R-7 was reported as met, full stop, on 270 observed rooms. It was met on the
occupancy regime it had been run in. Running the same experiment with every room at its design head
count fails \textbf{231 of 270}, and the shape of the failure is the result.

The two failing clauses are \textbf{disjoint}: 215 observed rooms were themselves \emph{accused of a
fault they did not have}, and 16 more sat outside the 5 \% airflow band at a worst deviation of
0.058919: 215 plus 16 is the 231. They attribute in opposite directions, and the supervisor-dead
pass is what separates them. The accusations run 85 of 90 rooms with the supervisor \emph{dead}
against 130 of 180 with it live, so they belong to the regime and the monitor rather than to the
response to the fault: compensation slightly \emph{suppresses} them. The airflow breaches run 0 of
90 dead against 16 of 180 live, so that half is the compensation's own cost and nothing else's.

\textbf{So the larger half of R-7's design-occupancy failure is R-3's false-alarm rate, measured by a
different experiment, in a regime R-3's own replay never visits.} That is the second time this report
has had to say it: §\ref{sec:11.2} records the same pattern inside R-2, where two of five design-occupancy
seeds fail because every room is latched by false alarms and the trim lever is gone. \textbf{R-3 is
not one failing requirement among four; it is the one whose failure the other three inherit}, and
stating it that way is worth more than three independently reported shortfalls. ADR-0009 predicted
exactly this collision and noted that no experiment asserted R-7 in that regime. One does now.

Two bounds on the finding. The first is that \texttt{CompensationPlan.worstHealthyFlowDeviation()}
\textbf{cannot falsify this requirement} and should stop being quoted beside the figure that can: it
reaches 0.0500 and no further in any regime, because the solver caps the fan at the speed that makes
the figure exactly the tolerance and then reports that cap. The second is that the 5 \% bound's
\emph{reference plant} is ambiguous: the solver bounds its claim over each latched room's
latch-time opening, the independent recomputation over its pre-latch one, and the two coincide until
a second room latches inside one episode, which the drawn weekday never does and design occupancy
routinely does. Both are defensible readings of ``before the compensation'' and ADR-0008 never says
which. That is an ADR-level question, not a code change.

\subsection{The third diagnostic signal is blind after a confirmation (\texorpdfstring{\S\ref{sec:13.6}}{13.6} in full)}\label{sec:I.5}


§\ref{sec:7.3} records that the stuck-damper leaf fires on two signals rather than three, because a probe can
only be requested once a case is latched and no step has been judged at classification time. The
obvious repair is to reorder: latch as usual, withhold the diagnosis, ask for the probe immediately,
and name the fault when the step lands. \textbf{It was built, measured over 176 paired runs on
identical seeds, and not kept.} Detection moves in 0 of 88 pairs and the fault is named on a
genuinely judged step in 28 of the 39 runs that defer, against none before. And the settled
classification is identical in all 88 pairs, because the judged step reads \texttt{MISSING} for a
stuck damper, a blocked filter, nine rooms whose dampers are in perfect health, and the run with
nothing injected.

The cause is a located defect older than the ordering, and it is the useful part.
\texttt{ZoneObserver.\allowbreak recordStepAmplitude} measures the amplitude's \emph{measured} leg against
\texttt{predictedAtStep} while the predicted leg is also against \texttt{predictedAtStep}, so the
measured leg carries the standing residual as a constant offset. And the anchor-gain pull that
would make that baseline legitimate is switched off while a case is open, so at the confirmation the
prediction is stale by exactly the residual that confirmed it. \texttt{ZoneObserver.measuredAtStep}
is assigned for this purpose and never read. On a 600 ppm offset with a healthy damper that answers
its probe by dropping 300 ppm, the amplitude ratio runs from $-7.0$ to $-0.8$.

\textbf{And correcting that baseline makes the reorder worse, not better}, which is the finding worth
defending. Measured as an unlanded counterfactual, the corrected signal takes T-15's diagonal from
18/44 to 15/44: it reads \texttt{MATCHED} on 3 of 8 genuinely stuck dampers and \texttt{MISSING} on 5
of 11 healthy ones, because by the time a probe can be commanded the supervisor has degraded the
zone, raised the fan and trimmed the other five, so the room's answer is mostly the compensation's. The 18/44 baseline reproduces on the current code; the 15/44 was measured on a one-line patch that was then reverted, so it is the one figure in this report that cannot be re-derived from the repository as it stands.
This is the masking the latching argument in §\ref{sec:13.8} already rests on, arriving one level down, in the
diagnosis rather than in the clearing policy. \textbf{A three-signal stuck-damper diagnosis therefore
needs a step judged \emph{before} a case opens} (the inner loop's own trims, read against a
baseline that survives them) rather than a probe requested after the response to the fault has
begun. That is design work, and it is the one concrete thing this report would hand to a successor.

\subsection{Scope, and the debts taken knowingly (\texorpdfstring{\S\ref{sec:13.9}}{13.9} in full)}\label{sec:I.6}

\textbf{Legitimate scope.} Six rooms on one floor behind one air handler is not a building, but the
shared fan makes R-7 non-trivial: keeping five rooms unchanged while compensating the sixth is
something the supervisor must actively do. Simulated time through \texttt{ClockPort} is legitimate
too, but the reproducibility it buys has to be stated precisely, because the obvious phrasing is
false of half the system. \textbf{In the offline harness the replay is reproducible bit for bit}:
the clock is a \texttt{SimulatedClock}, the bus is a list and the iteration order is fixed, so the
draws are identical run to run, which is exactly what R-3's false-alarm count needs.
\textbf{In the containerised demo profile it is not}, and that is measured rather than suspected:
four runs on identical \texttt{SEED}, \texttt{TIME\_SCALE} and \texttt{START\_AT} confirmed A115 at
simulated 10:16, 10:17, 10:19 and 10:23, a spread of seven simulated minutes. The cause is that
the sensor model draws once per Euler sub-step and the sub-step count follows \emph{real} elapsed
time, so host scheduling decides how many draws a simulated minute contains. The offline number is
the one every requirement is argued from; the demo profile is for showing the system work, not for
measuring it. What simulated time does \emph{not} exercise either way is anything that
appears only over real duration: scheduler jitter, clock drift, broker backpressure, memory growth.
R-6's cycle time is the only result in the plan that touches timing at all.

\textbf{Debts.} Skipping a DTO layer is justified by lock-step single-author deployment, but it makes
a domain refactor a wire-format break and is wrong the moment any consumer deploys independently.
Having no message replay is similarly contingent: R-3's evidence comes from a seeded harness rather
than from the broker, so the evaluation never needs replay, but in a real deployment a monitor
restart loses the residual history it was accumulating, and R-5 then rests entirely on the journal.
Three smaller debts, all three of which have since been \textbf{corrected} and are kept here because
what they were is worth more than that they are gone. The first: \texttt{Topology} could express a
synthetic zone set larger than six but nothing in any container wired it, so the capability existed
and R-6 could not be run from the deployed system. \texttt{ZONE\_COUNT} is now an environment
variable all \factAuthoredServices{} project-authored containers read, and the stack was brought up and watched at 12
and at 58 zones, 58 being a hard ceiling that belongs to BuildSim rather than to the monitor,
since it refuses to host equipment in a room its floorplan does not have. \textbf{R-6's breaking
point therefore remains a harness measurement}, and the two ceilings are three orders of magnitude
apart. The second: ADR-0005 and ADR-0006 argued in the superseded access-control vocabulary (doors,
badges, requirement IDs that no longer exist) against requirements the project had replaced. Both
were retargeted, and then checked line by line against the source tree, which found that the
retargeted prose made three further wrong statements about the code: an interface inventory listing
five where the tree has fifteen, a claim that five components carry the full
\texttt{boundary}/\texttt{control}/\texttt{entity} structure, and device-process line counts given as
``around a hundred each'' against 297 and 237. \textbf{Correcting a document can introduce fresh
errors of its own}, and only a pass against the tree finds them. The third, that the monitor's
parameter record documented its plausibility ceiling as 1.0 while the code used 1.25, sat in the one
parameter that decides how desensitised the detector is, which is the worst place to have two
different numbers written down: \texttt{MonitorParameters} now reads 1.25 in both its javadoc and its
defaults, and the same pass removed two further false claims from that record (ADR-0010).

\textbf{A fourth of the same shape was found and closed in this pass, and it was the most expensive
of them.} The probe step had \emph{two} numbers (\texttt{MonitorParameters.probeStepIncrease} at
0.20 and \texttt{SupervisionService.PROBE\_STEP} at 0.25), and the consequence was worse than
duplication, because the monitor \emph{asks} in the evidence an operator reads while the supervisor
\emph{acts}. The system published ``step to 0.35'' on \texttt{faults/\{zone\}}, journalled that, and
commanded 0.40: \textbf{it was telling the operator one thing and the building another}. There is one
owner now, in \texttt{shared}, because the amplitude has to mean the same thing to the component that
requests it and the one that commands it. Why the divergence survived is worth keeping: the request
is prose inside \texttt{FaultStateChanged.evidence}, a \texttt{List<String>}, so the supervisor could
not read the amplitude off the wire without a schema change four components decode. The change was
taken on the operator-facing argument and \textbf{not} on a measurement (no requirement verdict
moved), and it had one measured side effect worth more than the fix: a 0.05 change in one amplitude
moved \textbf{ten of T-34b's design-occupancy rows from pass to FAIL}, by causing one extra room to be
falsely accused on one seed. That is R-3's false-alarm rate again, and it says the design-occupancy
regime is not merely failing but \emph{unstable}.

\textbf{One debt was addressed, and what it cost to try is more interesting than the debt.}
T-15 found that \texttt{ZoneMonitor.stepShapeUnchanged()} collapsed ``no evidence'' into ``evidence
of a missing response'': it returned \texttt{false} when the step-amplitude ratio was \texttt{NaN},
and \texttt{NaN} is what the observer reports when \emph{no step has ever been judged}. The
classifier's stuck-damper leaf read that as the commanded step having produced no matching
response, so \texttt{STUCK\_DAMPER} could be emitted for a room the monitor never asked to move:
9 of 57 drift runs confirm by exactly that route.

The signal is now a three-valued type, so the conflation cannot be written by accident again. The
\emph{strict} fix, requiring positive evidence of an unanswered step, was then tried, and
\textbf{it emptied the leaf}: the one cell on T-15's diagonal went to \texttt{UNCLASSIFIED} and the
tree classified nothing at all. The cause is an ordering problem, not a logic one. A probe is
requested through the evidence on \texttt{FaultStateChanged}, so a case must be open before a step
is commanded, and \textbf{the step-amplitude signal therefore arrives after the decision it exists
to inform}. The \texttt{NaN} was not an incidental bug; it was the only thing making that leaf
reachable, and the taxonomy has been quietly running on two signals rather than three for its whole
life.

What is in the repository is the honest middle: the leaf still fires on localisation and flat fan
power, and every diagnosis reached without a judged step now carries an explicit entry saying it
rests on two signals, that the room has not been asked to move and so has not failed to answer, and
that a probe has been requested to settle it. Behaviour is unchanged (T-15's, T-24's and T-25's
result files reproduce byte for byte, so no number in this report moved), and the weakness is now
in the work order an operator reads instead of inside a boolean. \textbf{The real fix is to command
a probe step before classifying rather than after}, which reorders the diagnosis pipeline and would
trade detection delay for diagnostic confidence. That is a design change with its own experiment
attached, and it is the first thing this system should do next.
